\documentclass[12pt]{article}
\usepackage[utf8]{inputenc}
\usepackage[T1]{fontenc}
\usepackage[margin=1in]{geometry}
\usepackage{setspace}
\usepackage{xcolor}
\usepackage[colorlinks=true, urlcolor=blue!60!black, linkcolor=blue!60!black]{hyperref}
\setlength{\parindent}{0em}
\setlength{\parskip}{0.45em}
\renewcommand{\baselinestretch}{1.0}
\begin{document}
\begin{center}
{\Large \textit{United States v. Rahimi}, 602 U.S. 680 (2024)\par}
\vspace{0.3em}
{\normalsize Case Brief --- One-Page Edition\par}
\end{center}

\textbf{Memory Jogger.} \textit{Bruen} clarified: people a court finds dangerous to others may be temporarily disarmed (\href{https://www.law.cornell.edu/supremecourt/text/22-915}{602 U.S. 680}).

\textbf{Facts.} After assaulting his girlfriend, Zackey Rahimi became subject to a protective order finding he posed ``a credible threat'' to her safety; he kept guns and was linked to five shootings.

\textbf{Procedural History.} After \textit{Bruen}, the Fifth Circuit vacated his \S\ 922(g)(8) conviction.

\textbf{Votes.} Reversed, 8--1: Roberts for the Court; five concurrences; Thomas dissented.

\textbf{Holding.} Section 922(g)(8)(C)(i) is valid as applied and survives a facial challenge; laws need fit historical principles, not a ``historical twin.''

\textbf{Reasoning.} The precedents ``were not meant to suggest a law trapped in amber.'' Surety and going-armed laws show that those who threaten violence may be disarmed, and the statute matches their ``why'' and ``how.'' Barrett: ``[h]istorical regulations reveal a principle, not a mold''; Jackson saw ``little method to \textit{Bruen}'s madness.'' Thomas found no analogue imposing a comparable burden: surety laws left arms intact, and affray laws punished public terror, not domestic conduct.

\textbf{Examples.} \emph{Same-Side:} red-flag orders after a hearing. \emph{Opposite-Side:} disarming a nonviolent felon. \emph{Fence-Sitter:} orders without a threat finding.

\textbf{Critique.} Originalists call it ``\textit{Bruen}-lite'' and see many originalisms on the Court; progressives such as Blocher, Ruben, and Siegel say the level-of-generality problem remains and that Founding-era analogues privilege an era that excluded women.

\textbf{Key Quotation.} ``Taken together, the surety and going armed laws confirm what common sense suggests: When an individual poses a clear threat of physical violence to another, the threatening individual may be disarmed.'' --- Chief Justice Roberts, Opinion of the Court, slip op.\ at 13.

\end{document}
